% Methodology

\section{Overview}

This chapter describes the simulation model architecture, agent behaviors, experimental design, metrics, and statistical analysis methods employed in this research. We implement a grid-based multi-agent fruit harvesting simulation using the Mesa 3.0 framework, conduct a comprehensive factorial experiment across multiple parameter dimensions, and analyze results using rigorous statistical techniques.

\section{Simulation Model Architecture}

\subsection{Environment}

The simulation environment consists of a two-dimensional square grid with the following characteristics:

\begin{itemize}
    \item \textbf{Grid size:} Parameterizable dimensions (default 50×50 cells, scalable from 10×10 to 200×200)
    \item \textbf{Topology:} Non-toroidal (agents cannot wrap around edges)
    \item \textbf{Cell occupancy:} Multiple agents and one fruit resource per cell (MultiGrid)
    \item \textbf{Obstacles:} None in baseline configuration
\end{itemize}

\subsection{Fruit Resources}

Fruit resources occupy grid cells and exhibit two dynamics modes:

\subsubsection{Static Dynamics}

In Static mode, fruit resources are placed at initialization and do not regenerate after harvest. This represents one-time harvest scenarios such as grain fields. Key parameters:

\begin{itemize}
    \item \textbf{Initial fruit density:} Fraction of cells containing fruit at $t=0$ (values: 0.10, 0.20, 0.30)
    \item \textbf{Regeneration:} None ($p_{regen} = 0.0$)
\end{itemize}

\subsubsection{Replenishing Dynamics}

In Replenishing mode, harvested fruit can regenerate with a specified probability per time step. This represents continuous ripening scenarios such as strawberry fields. Key parameters:

\begin{itemize}
    \item \textbf{Initial fruit density:} Same as Static mode
    \item \textbf{Regeneration probability:} $p_{regen} \in \{0.01, 0.05, 0.10\}$ per step for each empty fruit cell
\end{itemize}

\subsection{Harvester Agents}

Harvester agents are autonomous entities that navigate the grid, harvest fruit, and communicate with teammates. Agent behaviors follow Mesa 3.0 conventions (no mandatory \texttt{unique\_id}, activation via \texttt{model.agents.shuffle\_do('step')}).

\subsubsection{Movement}

Agents move one cell per time step using Moore neighborhood (8-directional movement). Movement strategy:

\begin{enumerate}
    \item If agent has a target position (received via communication), move toward target using greedy strategy (reduce Manhattan distance).
    \item Otherwise, perform random walk to adjacent cells.
\end{enumerate}

\subsubsection{Harvesting}

At each time step, agents check their current cell for available fruit. If fruit is present and available, the agent harvests it (fruit becomes unavailable, agent's harvest count increments).

\subsubsection{Communication}

Agents communicate with teammates within a specified communication range using Chebyshev distance (maximum of absolute coordinate differences). Communication protocol:

\begin{enumerate}
    \item Agent identifies all teammates within communication range $r_{comm}$ (values: 0--8 cells).
    \item Agent searches for nearest available fruit within perception range.
    \item If fruit is found, agent broadcasts fruit location to all neighbors without current targets.
    \item Receiving agents update their target position and navigate toward the communicated fruit location.
\end{enumerate}

Message tracking: Each agent maintains counters for messages sent and received per step, enabling communication overhead analysis.

\section{Experimental Design}

\subsection{Independent Variables}

The experiment employs a full factorial design with the following independent variables:

\begin{table}[h]
\centering
\caption{Independent Variables and Levels}
\begin{tabular}{lll}
\toprule
\textbf{Variable} & \textbf{Levels} & \textbf{Description} \\
\midrule
Communication range & 0, 1, 2, 3, 4, 5, 6, 7, 8 & Cells (Chebyshev distance) \\
Team size & 10, 20, 30 & Number of harvester agents \\
Resource dynamics & Static, Replenishing & Regeneration mode \\
Regeneration prob. & 0.00, 0.01, 0.05, 0.10 & Per step, per empty cell \\
Fruit density & 0.10, 0.20, 0.30 & Fraction of cells with fruit \\
\bottomrule
\end{tabular}
\end{table}

\textbf{Note:} Regeneration probability applies only to Replenishing dynamics. For Static dynamics, $p_{regen} = 0.0$ regardless of specified value.

\subsection{Design Structure}

\begin{itemize}
    \item \textbf{Static condition:} $9 \text{ (range)} \times 3 \text{ (agents)} \times 3 \text{ (density)} \times 30 \text{ (seeds)} = 2,430$ runs
    \item \textbf{Replenishing condition:} $9 \times 3 \times 3 \times 4 \text{ (regen)} \times 30 = 9,720$ runs
    \item \textbf{Total:} 12,150 runs
\end{itemize}

\subsection{Pilot Study}

A pilot study with 5 seeds per configuration is conducted first to:
\begin{enumerate}
    \item Validate metrics collection and data pipeline
    \item Estimate effect sizes for power analysis
    \item Assess computational runtime and identify bottlenecks
    \item Refine parameter ranges if necessary
\end{enumerate}

\subsection{Simulation Parameters}

\begin{table}[h]
\centering
\caption{Fixed Simulation Parameters}
\begin{tabular}{ll}
\toprule
\textbf{Parameter} & \textbf{Value} \\
\midrule
Grid width & 50 cells \\
Grid height & 50 cells \\
Time steps & 500 \\
Agent perception range & Unlimited (can see all fruit) \\
Movement speed & 1 cell per step \\
Harvest capacity & 1 fruit per step \\
\bottomrule
\end{tabular}
\end{table}

\section{Metrics and Data Collection}

\subsection{Per-Step Metrics}

The following metrics are collected at each time step using Mesa's \texttt{DataCollector}:

\begin{itemize}
    \item \textbf{total\_yield:} Cumulative fruit harvested by all agents
    \item \textbf{messages\_this\_step:} Messages sent this step
    \item \textbf{cumulative\_messages:} Total messages sent since $t=0$
    \item \textbf{remaining\_fruit:} Available fruit count
    \item \textbf{harvested\_cells\_this\_step:} Cells harvested this step
    \item \textbf{largest\_component\_ratio:} $|V_{max}| / N$ where $V_{max}$ is the largest connected component in the communication graph
    \item \textbf{avg\_hops:} Mean shortest-path length for message delivery
    \item \textbf{mean\_agent\_load:} Mean harvest count per agent
    \item \textbf{gini\_yield:} Gini coefficient of harvest distribution
    \item \textbf{avg\_travel\_distance:} Mean movement steps per agent
\end{itemize}

\subsection{Per-Run Summary Metrics}

Aggregated metrics computed after each run:

\begin{itemize}
    \item \textbf{final\_yield:} Total yield at final time step
    \item \textbf{steady\_state\_yield:} Mean yield over last 25\% of steps
    \item \textbf{time\_to\_depletion:} Step when remaining fruit reaches 0 (Static only)
    \item \textbf{mean\_messages\_per\_step:} Average messages per step
    \item \textbf{yield\_per\_message:} $\text{final\_yield} / \max(\text{cumulative\_messages}, 1)$
    \item \textbf{peak\_component\_ratio:} Maximum connectivity achieved
\end{itemize}

\subsection{Metric Formulas}

\subsubsection{Largest Component Ratio}

Communication graph $G = (V, E)$ where $V$ is the set of agents and $(i, j) \in E$ if $d_{Chebyshev}(pos_i, pos_j) \leq r_{comm}$. Largest component ratio:

\begin{equation}
\text{LCR} = \frac{|V_{max}|}{|V|}
\end{equation}

where $V_{max}$ is the largest connected component found via depth-first search.

\subsubsection{Gini Coefficient}

For harvest counts $h_1 \leq h_2 \leq \cdots \leq h_n$:

\begin{equation}
G = \frac{2 \sum_{i=1}^{n} i \cdot h_i}{n \sum_{i=1}^{n} h_i} - \frac{n+1}{n}
\end{equation}

Values range from 0 (perfect equality) to 1 (maximum inequality).

\section{Statistical Analysis}

\subsection{Primary Outcome}

The primary outcome variable is \textbf{final\_yield}. Secondary outcomes include \textbf{yield\_per\_message}, \textbf{largest\_component\_ratio}, and \textbf{steady\_state\_yield}.

\subsection{Analysis Pipeline}

\subsubsection{Descriptive Statistics}

\begin{itemize}
    \item Mean, standard deviation, median, and interquartile range for each configuration
    \item Visualization: line plots (yield vs. range), heatmaps (yield over range × agents), time-series (yield over time)
\end{itemize}

\subsubsection{Inferential Statistics}

\textbf{ANOVA:} Multi-way analysis of variance to test main effects and interactions:

\begin{itemize}
    \item \textbf{Static:} $\text{final\_yield} \sim \text{comm\_range} \times \text{num\_agents} \times \text{fruit\_density}$
    \item \textbf{Replenishing:} Add $\text{regen\_prob}$ as factor
\end{itemize}

\textbf{Mixed-Effects Model:} Linear mixed model with random intercept for seed:

\begin{equation}
Y_{ijklm} = \mu + \alpha_i + \beta_j + \gamma_k + \delta_l + (\alpha\beta)_{ij} + \cdots + u_m + \epsilon_{ijklm}
\end{equation}

where $u_m \sim N(0, \sigma_u^2)$ is the random effect for seed $m$.

\textbf{Post-Hoc Tests:} Tukey HSD for pairwise comparisons of significant factors.

\textbf{Significance Level:} $\alpha = 0.05$ with 95\% confidence intervals.

\textbf{Effect Sizes:} Partial $\eta^2$ for ANOVA, Cohen's $d$ for pairwise comparisons.

\subsubsection{Assumption Checking}

\begin{itemize}
    \item \textbf{Normality:} Shapiro-Wilk test on residuals; Q-Q plots
    \item \textbf{Homoscedasticity:} Levene's test; residual plots
    \item \textbf{Independence:} Verified by experimental design (separate seeds)
\end{itemize}

If assumptions are violated, apply transformations (log, square root) or use non-parametric alternatives (Kruskal-Wallis, Dunn's test with Holm correction).

\section{Implementation Details}

\subsection{Software and Tools}

\begin{itemize}
    \item \textbf{Simulation:} Mesa 3.0 (Python)
    \item \textbf{Data analysis:} pandas, NumPy, SciPy, statsmodels
    \item \textbf{Visualization:} Matplotlib, Seaborn, Solara
    \item \textbf{Statistical testing:} SciPy, statsmodels
    \item \textbf{Version control:} Git, GitHub
\end{itemize}

\subsection{Computational Resources}

\begin{itemize}
    \item \textbf{Hardware:} Apple M-series laptop
    \item \textbf{Parallelization:} Multi-core execution for independent runs
    \item \textbf{Estimated runtime:} 1.7--6.7 hours for 10,000 runs (0.5--2.0 seconds per run)
\end{itemize}

\subsection{Reproducibility}

\begin{itemize}
    \item Fixed random seeds for all runs
    \item Configuration files (JSON) stored with results
    \item Metadata logging (timestamp, parameters, software versions)
    \item Source code archived in version control repository
\end{itemize}

\section{Ethical Considerations}

This research involves computational simulations only, with no human subjects or animal testing. All software used is open-source or freely available for academic use. Results and code will be made publicly available to support reproducibility and future research.

